\documentclass[10.5pt,a4paper]{article}

\usepackage[utf8]{inputenc}
\usepackage[T1]{fontenc}
\usepackage{geometry}
\usepackage{booktabs}
\usepackage{xcolor}
\usepackage{colortbl}
\usepackage{fancyhdr}
\usepackage{titlesec}
\usepackage{hyperref}
\usepackage{tcolorbox}
\usepackage{enumitem}
\usepackage{array}
\usepackage{parskip}
\usepackage{amsmath}
\usepackage{amssymb}
\usepackage{mathtools}

\geometry{left=2.6cm,right=2.6cm,top=2.4cm,bottom=2.4cm}

\definecolor{primary}{RGB}{15,23,42}
\definecolor{accent}{RGB}{30,64,175}
\definecolor{muted}{RGB}{100,116,139}
\definecolor{lightbg}{RGB}{248,250,252}
\definecolor{mathbg}{RGB}{241,245,249}
\definecolor{bordercol}{RGB}{203,213,225}
\definecolor{greenbg}{RGB}{240,253,244}
\definecolor{greenborder}{RGB}{134,239,172}
\definecolor{greentext}{RGB}{22,101,52}
\definecolor{redbg}{RGB}{255,241,242}
\definecolor{redborder}{RGB}{252,165,165}
\definecolor{redtext}{RGB}{153,27,27}

\hypersetup{colorlinks=false}

\titleformat{\section}
  {\large\bfseries\color{primary}}{}{0em}{}[{\color{accent}\vspace{2pt}\hrule height 0.7pt}]
\titlespacing{\section}{0pt}{1.4em}{0.5em}

\titleformat{\subsection}
  {\normalsize\bfseries\color{accent}}{}{0em}{}
\titlespacing{\subsection}{0pt}{0.9em}{0.2em}

\pagestyle{fancy}
\fancyhf{}
\fancyhead[L]{\small\color{muted}\textit{Confidencial}}
\fancyhead[R]{\small\color{muted}Junio 2026}
\fancyfoot[C]{\small\color{muted}\thepage}
\renewcommand{\headrulewidth}{0pt}

\tcbuselibrary{skins,breakable}

\newtcolorbox{mathbox}{
  colback=mathbg, colframe=bordercol, boxrule=0.6pt,
  arc=2pt, left=8pt, right=8pt, top=6pt, bottom=6pt, breakable}

\newtcolorbox{greenbox}{
  colback=greenbg, colframe=greenborder, boxrule=0.6pt,
  arc=2pt, left=8pt, right=8pt, top=6pt, bottom=6pt, breakable}

\newtcolorbox{redbox}{
  colback=redbg, colframe=redborder, boxrule=0.6pt,
  arc=2pt, left=8pt, right=8pt, top=6pt, bottom=6pt, breakable}

\newtcolorbox{notebox}{
  colback=lightbg, colframe=bordercol, boxrule=0.6pt,
  arc=2pt, left=8pt, right=8pt, top=6pt, bottom=6pt, breakable}

\setlength{\parskip}{0.38em}
\setlength{\parindent}{0em}
\renewcommand{\arraystretch}{1.3}
\setlist[itemize]{nosep, topsep=1pt, leftmargin=1.5em, label={\color{accent}\textbullet}}

% ============================================================
\begin{document}

% PORTADA
\thispagestyle{empty}
\vspace*{1.1cm}

{\color{muted}\small\textbf{CONFIDENCIAL}\quad|\quad Junio 2026}

\vspace{0.9cm}

{\fontsize{28}{33}\selectfont\bfseries\color{primary}
Mantenimiento predictivo\\de flota vehicular}

\vspace{0.4cm}

{\fontsize{13}{18}\selectfont\color{accent}
Fundamentos fisicos, matematicos y caso economico}

\vspace{0.5cm}
{\color{bordercol}\hrule height 1pt}
\vspace{0.5cm}

{\large\color{muted}
Un sistema que modela como se degrada cada pieza en tiempo real,\\
cuantifica la vida residual con incertidumbre, y decide\\
el momento optimo de intervencion -- con matematicas derivadas\\
desde primeros principios y un caso economico trazable.
}

\vspace{1.0cm}

\begin{notebox}
\textbf{Proposito de este documento.}
No hay pitch aqui. Lo que hay es la matematica real que subyace a lo que construimos
-- derivaciones desde primeros principios, no cajas negras --, y el calculo riguroso
de que significa economicamente adoptar un esquema de mantenimiento predictivo
sobre una flota. Para quien tiene ingenieros excelentes: aqui esta todo lo que
necesitan para evaluar si el enfoque es correcto.
\end{notebox}

\vspace{0.9cm}

\begin{tabular}{@{}p{0.6cm}p{12.8cm}@{}}
  \color{accent}\textbf{I}   & El costo real de no hacer nada: la aritmetica del mantenimiento reactivo \\[4pt]
  \color{accent}\textbf{II}  & Por que el calendario fijo no es la solucion: el modelo correcto de falla \\[4pt]
  \color{accent}\textbf{III} & Distribucion de vida desde primeros principios (Weibull y teorema de extremos) \\[4pt]
  \color{accent}\textbf{IV}  & Fisica de falla componente a componente \\[4pt]
  \color{accent}\textbf{V}   & Inferencia con datos censurados y regresion de supervivencia \\[4pt]
  \color{accent}\textbf{VI}  & Vida util remanente como distribucion: procesos de Wiener y Gamma \\[4pt]
  \color{accent}\textbf{VII} & Estimacion online bajo mantenimiento reactivo: el regimen de arranque \\[4pt]
  \color{accent}\textbf{VIII}& Modelo jerarquico bayesiano de flota multimarca \\[4pt]
  \color{accent}\textbf{IX}  & Decision optima: cuando intervenir y cuanto cuesta no hacerlo \\[4pt]
  \color{accent}\textbf{X}   & Cuantificacion del ahorro esperado \\[4pt]
  \color{accent}\textbf{XI}  & Por que este enfoque supera a los existentes \\
\end{tabular}

\newpage

% ============================================================
\section{El costo real de no hacer nada}

Antes de cualquier ecuacion, el problema economico. Los benchmarks que siguen
provienen del Departamento de Energia de EE.UU. (US DOE), McKinsey \& Company (2020),
el Electric Power Research Institute (EPRI) y estudios de operadores de flota
documentados entre 2022 y 2025. Se citan con precision porque son las mismas
fuentes que usan los evaluadores institucionales.

\subsection{La estructura de costos del mantenimiento reactivo}

Una reparacion de emergencia no cuesta lo mismo que la misma reparacion programada.
Los componentes de costo son acumulativos:

\begin{center}
\small
\begin{tabular}{@{}p{4.8cm}p{5.8cm}r@{}}
\toprule
\textbf{Componente} & \textbf{Mecanismo} & \textbf{Multiplicador tipico} \\
\midrule
Mano de obra & Tiempo extra, servicio nocturno o de fin de semana & $1.5\times$--$2.0\times$ \\
Refacciones & Envio urgente, sin poder negociar precio & $1.2\times$--$1.4\times$ \\
Dano secundario & Falla en cascada al componente adyacente & $3\times$--$5\times$ \\
Tiempo muerto & Unidad varada, perdida de ingreso por viaje & cuantificable \\
Grua y arrastre & Costo medio por evento en ruta & fijo por evento \\
\bottomrule
\end{tabular}
\end{center}

\medskip
El benchmark mas citado y replicado: \textbf{una reparacion reactiva cuesta en
promedio 4.8 veces mas que la misma reparacion ejecutada como mantenimiento
programado} (US DOE; replicado por Deloitte y operadores de flota de clase 4--8).
El US DOE documenta adicionalmente:

\begin{redbox}
\begin{itemize}
  \item Reduccion de costos de mantenimiento: \textbf{25\%--30\%} al pasar de reactivo a predictivo.
  \item Eliminacion de fallas no planeadas: \textbf{70\%--75\%}.
  \item Reduccion de tiempo muerto no planeado: \textbf{35\%--45\%}.
  \item Incremento en produccion/disponibilidad: \textbf{20\%--25\%}.
  \item ROI documentado: \textbf{10:1} en implementaciones bien ejecutadas.
\end{itemize}
\vspace{2pt}
{\small\color{muted} Fuente: US DOE, \textit{Operations \& Maintenance Best Practices Guide}, ediciones 2.0--3.0.
Confirmado independientemente por EPRI (30\% de reduccion sobre PM periodico;
50\% sobre reactivo puro) y por McKinsey \& Company (2020): reduccion
de costos de mantenimiento 10\%--40\%, tiempo muerto 50\%--70\%.}
\end{redbox}

\subsection{El multiplicador de dano secundario: la pieza mas subestimada}

Una falla de DPF en ruta no solo cuesta el sensor de NOx o el filtro.
Cuesta el derate que deja la unidad a baja potencia, la grua, el tiempo muerto,
el costo de oportunidad del viaje no completado, y con frecuencia dano al
turbocompresor por contaminacion del aceite. El costo del componente primario
puede ser la fraccion mas pequena del total. El multiplicador $3\times$--$5\times$
de dano secundario es la razon de que las fallas en ruta de flota disel sean
desproporcionadamente caras.

% ============================================================
\section{Por que el calendario fijo no es la solucion}

Existe una alternativa inmediata al mantenimiento reactivo: el mantenimiento
preventivo por calendario (cambio cada $N$ km o $M$ meses). Es mejor que
nada, pero tiene dos fallas estructurales.

\textbf{Falla 1: remplaza piezas con vida util remanente.}
IBM y operadores de flota documentan que el 30\% de las intervenciones
preventivas sustituyen componentes que aun tienen mas del 40\% de su vida til.
Eso no es mantenimiento: es desperdicio de refacciones.

\textbf{Falla 2: no distingue contexto operativo.}
Un km en ruta de montana con carga completa no equivale a un km en carretera plana.
Una balata sometida a 200 frenadas de 80 a 0 km/h tiene mas dano acumulado
que una con el doble de km pero condiciones suaves. El calendario no sabe eso.

\medskip
La unica solucion que resuelve ambas fallas es medir el estado real de cada
componente y actualizar el pronostico de vida con cada nueva observacion.
Eso es lo que construimos.

% ============================================================
\section{El modelo correcto de falla: la fisica como punto de partida}

Sea $T \geq 0$ el tiempo (o km, u horas-motor) hasta la falla de un componente.
Lo que importa fisicamente no es $T$ directamente, sino la \textbf{tasa de riesgo}:

\begin{mathbox}
$$h(t) = \lim_{\Delta t \to 0} \frac{P(t < T \leq t+\Delta t \mid T > t)}{\Delta t}
       = \frac{f(t)}{R(t)} = -\frac{d}{dt}\ln R(t)$$

Dado que el componente llego vivo al tiempo $t$, $h(t)$ es la tasa a la que esta
fallando \textit{ahora mismo}. Especificar $h(t)$ equivale a especificar el modelo completo:
$$R(t) = \exp\!\left(-\int_0^t h(u)\,du\right).$$
\end{mathbox}

La \textbf{vida util remanente} (RUL) -- lo que el operador necesita -- es:

\begin{mathbox}
$$\mathrm{RUL}(t) = \mathbb{E}[T - t \mid T > t]
                 = \frac{\displaystyle\int_t^\infty R(u)\,du}{R(t)}$$

Esto \textit{no} es $\mathrm{MTTF} - t$. Para riesgo creciente ($\beta > 1$, desgaste), la RUL
decrece con la edad. Para riesgo constante (falla aleatoria, $\beta = 1$), la pieza
no envejece y el mantenimiento por calendario no ayuda en absoluto.
\end{mathbox}

\subsection{Por que Weibull es la distribucion inevitable: valores extremos}

Un componente real es una coleccion de $N$ micro-eslabones (granos, asperezas,
volumenes elementales). Falla cuando el mas debil falla:
$T = \min(X_1, \dots, X_N)$. Si la cola izquierda de cada eslabon es ley de
potencia, $F_X(t) \approx (t/t_0)^\beta$:

\begin{mathbox}
$$R_T(t) = [1 - (t/t_0)^\beta]^N \;\xrightarrow{N \gg 1}\;
  \exp\!\left[-\left(\frac{t}{\eta}\right)^\beta\right], \qquad
  \eta = t_0\,N^{-1/\beta}$$

Teorema de Fisher--Tippett--Gnedenko (1928, 1943): la distribucion de minimos de
variables acotadas por abajo converge a Weibull. No hay eleccion; hay derivacion.
$\beta$ selecciona el regimen: $<1$ mortalidad infantil, $=1$ fallas aleatorias,
$>1$ desgaste. $\eta$ es la vida caracteristica.
\end{mathbox}

Consecuencia fisica: \textbf{efecto de tamano.} $\eta \propto N^{-1/\beta}$.
Piezas mas grandes (mas eslabones) son intrisecamente mas debiles. Esto se mide
en materiales -- es la razon de que la resistencia a la fractura tenga dispersion
Weibull (Weibull lo demostro estudiando resistencia de materiales, 1939--1951).

La vida residual esperada de un componente con edad $t$ tiene forma cerrada:
\begin{mathbox}
$$\mathbb{E}[T \mid T>t] = \eta\, e^{(t/\eta)^\beta}\,
  \Gamma\!\left(1+\tfrac{1}{\beta},\; (t/\eta)^\beta\right), \qquad
  \mathrm{RUL}(t) = \mathbb{E}[T\mid T>t] - t$$

$\Gamma(a,x) = \int_x^\infty w^{a-1}e^{-w}dw$ es la funcion gamma incompleta
superior. Esta es la expresion que el motor de pronostico evalua en linea
conforme cada componente envejece.
\end{mathbox}

% ============================================================
\section{Fisica de falla componente a componente}

El enfoque Physics-of-Failure (PoF) modela el mecanismo de degradacion desde
primeros principios. Su ventaja frente a modelos puramente data-driven: generaliza
a marcas y modelos sin necesitar millones de fallas por cada uno, porque la fisica
de Archard, Miner y Lundberg--Palmgren es la misma para cualquier freno de disco.

\subsection{Frenos -- Ley de Archard derivada del contacto de asperezas}

Dos superficies se tocan solo en las cimas de sus rugosidades. Bajo carga $W$,
las asperezas se deforman plasticamente hasta que la presion iguala la dureza
$H$. El area real de contacto es $A_r = W/H$, independiente del area aparente
(resultado de Bowden--Tabor, 1950). Al deslizar, cada junta de radio $a$ se rompe
en distancia $\sim 2a$; con probabilidad $\kappa$ arranca una particula hemisferica
de volumen $\frac{2}{3}\pi a^3$ (criterio de Rabinowicz):

\begin{mathbox}
$$\boxed{V = K\,\frac{W\,L}{H}}, \qquad K \equiv \frac{\kappa}{3}, \quad
  \frac{dh_w}{dt} = k\,p\,v = \frac{K}{H}\,p\,v$$

$h_w$ = profundidad desgastada, $p$ = presion de contacto, $v$ = velocidad de deslizamiento.
\end{mathbox}

La energia friccional es $\dot{e} = \mu_f p v$, por lo que el desgaste de balata
es proporcional a la energia disipada. Cada frenada contribuye:
$$\Delta h_w \propto \frac{E_{\mathrm{frenada}}}{H}, \qquad
  E_{\mathrm{frenada}} \approx \tfrac{1}{2}m\,(v_i^2 - v_f^2).$$

Se mide $v$ de GPS y J1939/OBD, desaceleracion del IMU, masa de los sensores de
eje (SPN 178/179). El acumulado de energia de frenado por unidad es el predictor
fisico correcto, no el odometro.

\subsection{Fatiga estructural -- Palmgren--Miner y conteo rainflow}

Para chasis y suspension, el IMU es el instrumento primario. El dano acumulado
(regla de Miner):
$$D = \sum_i \frac{n_i}{N_i} \geq 1 \;\Rightarrow\; \text{falla}, \qquad
  \sigma_a = \sigma_f'\,(2N_f)^b \quad\text{(Basquin, 1910)}$$

El \textbf{conteo rainflow} (Matsuishi--Endo, 1968; norma ASTM E1049) descompone
la senal del IMU en lazos cerrados de histeresis. El dano acumulado:
$$D_\beta = \alpha\sum_i h_i^\beta, \quad \beta \approx 3\text{ (grieta)},\;
  \beta \approx 5 \text{ (iniciacion)}.$$

El rainflow se calcula en el borde; se emiten histogramas de ciclos, no senal cruda.

\subsection{Rodamientos -- Lundberg--Palmgren es Weibull de eslabon mas debil}

La falla por fatiga de contacto se inicia bajo la superficie donde el cortante
ortogonal $\tau_0$ es maximo. Aplicando el argumento de eslabon mas debil al
volumen $V$:
$$\ln\frac{1}{S} \propto \frac{\tau_0^c N^e V}{z_0^h} \;\Longrightarrow\;
  \boxed{L_{10} = \left(\frac{C}{P}\right)^p,\quad p=3\text{ bolas},\;
  p=\tfrac{10}{3}\text{ rodillos}}$$

$C$ es dato del fabricante; $P$ se calcula de las cargas medidas por J1939.
Esta es la misma derivacion de la Parte~III aplicada al volumen subsuperficial.

\subsection{Aftertreatment (DPF/DEF/SCR) -- el quick win totalmente alimentable}

J1939 difunde sin pollear: $\Delta P$ del DPF (SPN~3251), nivel de DEF (SPN~1761),
derate por NOx (SPN~5246), presion y temperatura de aceite (SPN~100, 175),
conteo de regeneraciones. La tendencia de $\Delta P$ y un codigo pendiente de NOx
son senales predictivas dias antes del derate. \textbf{Este es el primer modulo:
dato totalmente alimentable hoy, sin calibracion adicional, atacando las fallas
mas caras del diesel moderno.}

% ============================================================
\section{Inferencia con datos censurados}

En una flota viva, la mayoria de los componentes no ha fallado al observarlos:
$T_i > t_i$. Esta es \textbf{censura por la derecha}. Ignorarla sesga la estimacion
hacia vidas cortas. La verosimilitud correcta es:

\begin{mathbox}
$$\mathcal{L}(\theta) = \prod_i h(t_i;\theta)^{\delta_i}\,R(t_i;\theta)$$

$\delta_i = 1$ si fallo, $\delta_i = 0$ si sigue vivo. Cada observacion aporta
$R(t_i)$; las que fallaron aportan ademas $h(t_i)$ en el instante exacto de falla.
\end{mathbox}

El modelo de \textbf{riesgos proporcionales de Cox} incorpora telemetria:
$$h(t \mid \mathbf{x}) = h_0(t)\,\exp(\boldsymbol\beta^\top\mathbf{x})$$

En cada tiempo de falla $t_{(j)}$, con conjunto en riesgo $\mathcal{R}_j$:
$$P(\text{falla}\;i \mid t_{(j)}) =
  \frac{e^{\boldsymbol\beta^\top\mathbf{x}_i}}
       {\sum_{k\in\mathcal{R}_j}e^{\boldsymbol\beta^\top\mathbf{x}_k}}$$

El riesgo base $h_0(t)$ se cancela exactamente. Se estima $\boldsymbol\beta$ sin
asumir ninguna forma funcional para $h_0$. El vector $\mathbf{x}$ incluye: energia
de frenado acumulada (Archard), dano rainflow--Miner, severidad de ruta (GPS),
carga media (J1939), marca y modelo. La fisica reduce la dimensionalidad y
la cantidad de datos necesaria.

% ============================================================
\section{Vida util remanente como distribucion: procesos de Wiener y Gamma}

Cuando se mide una variable que degrada (espesor de balata, $\Delta P$ del DPF,
capacidad de bateria), el pronostico pasa de poblacional a individual.

\textbf{Proceso de Wiener.} Degradacion lineal con ruido:
$X(t) = x_0 + \mu t + \sigma B(t)$. La falla ocurre al primer cruce del umbral
$a$. Por el principio de reflexion del browniano con deriva:

\begin{mathbox}
$$f_{T_a}(t) = \frac{(a-x_0)}{\sigma\sqrt{2\pi t^3}}
  \exp\!\left[-\frac{(a-x_0-\mu t)^2}{2\sigma^2 t}\right]
  \;\sim\; \mathrm{IG}\!\left(\frac{a-x_0}{\mu},\;\frac{(a-x_0)^2}{\sigma^2}\right)$$

La RUL tiene distribucion Gaussiana Inversa. No es un numero puntual: es una
distribucion completa, que da intervalos de confianza sobre el pronostico.
\end{mathbox}

La deriva $\mu$ no es un parametro libre: la fija la fisica. En frenos,
$\mu = k\langle pv\rangle$ (Archard). En el DPF, $\mu$ es la tasa de acumulacion
de hollin. La fisica provee el prior; los datos ajustan el residuo y la
incertidumbre $\sigma$.

\textbf{Proceso Gamma.} Para desgaste estrictamente monotono: $\Delta X \sim
\mathrm{Gamma}(\alpha\Delta t,\lambda)$. La confiabilidad en $t$ es la funcion
gamma incompleta regularizada evaluada en el umbral.

% ============================================================
\section{Estimacion online bajo mantenimiento reactivo: el regimen de arranque}

Esta es la situacion real durante los primeros meses de operacion del sistema:
la flota opera principalmente en modo reactivo, los datos de falla son escasos
y los datos censurados dominan. El sistema no espera a tener ``suficientes'' datos;
\textbf{aprende en tiempo real con lo que llega.}

\subsection{El filtro bayesiano como ecuacion maestra}

La densidad de creencia $p(x,t)$ sobre el estado de salud $x$ obedece la
\textbf{ecuacion de Fokker--Planck}:

\begin{mathbox}
$$\frac{\partial p}{\partial t} = -\frac{\partial}{\partial x}[\mu\,p]
  + \frac{1}{2}\frac{\partial^2}{\partial x^2}[\sigma^2 p]$$

Cada nueva observacion $y$ actualiza la densidad por Bayes puntual:
$p(x_t\mid y_{1:t}) \propto p(y_t\mid x_t)\,p(x_t\mid y_{1:t-1})$.

Propagar la densidad con Fokker--Planck y corregirla con Bayes en cada observacion
es exactamente el filtro de particulas.
\end{mathbox}

\subsection{El factor de reparacion: cuanto restaura cada intervencion}

No toda reparacion devuelve el componente a nuevo. La \textbf{edad virtual de
Kijima} modela esto: tras la $i$-esima intervencion,
$v_i = v_{i-1} + q\,x_i$, con $x_i$ el tiempo del ultimo ciclo. El parametro $q$
es el \textbf{factor de restauracion}: $q=0$ es reparacion perfecta,
$q=1$ es minima (el componente queda ``tan viejo como antes''),
$q>1$ es peor que antes (reparacion que dana).

\begin{notebox}
En el regimen de arranque con mantenimiento reactivo, cada reparacion es
una observacion de alta calidad: da el tiempo exacto de falla y el tipo de
intervencion. Esas observaciones son oro para calibrar $q$ por taller y por
tipo de componente. El sistema en modo reactivo \textbf{esta acumulando el
dataset mas valioso posible}: fallas reales con tiempos exactos. Cuando el
modelo tiene suficiente masa posterior, las alertas predictivas comienzan a
desplazar a las reactivas. La transicion es gradual y medible.
\end{notebox}

\subsection{Arquitectura de capas: ninguna recalcula todo}

\begin{center}
\small
\begin{tabular}{@{}p{3.6cm}p{3.5cm}p{3.5cm}p{2.2cm}@{}}
\toprule
\textbf{Capa} & \textbf{Que estima} & \textbf{Metodo} & \textbf{Cadencia} \\
\midrule
Filtrado de estado & RUL de cada unidad & KF / UKF / particulas & por evento \\
Eventos recurrentes & factor de reparacion $q$ & GRP / Kijima & por reparacion \\
Hiperparametros de flota & $\beta,\eta,\tau_\bullet,\boldsymbol\gamma$ & HMC / SG-MCMC & nocturno \\
Despliegue & posterior unidad nueva & red amortizada (SBI) & instantaneo \\
Politica & accion optima & POMDP / RL & por decision \\
\bottomrule
\end{tabular}
\end{center}

El filtrado es $O(1)$ por paso; los hiperparametros se refrescan en segundo
plano; la inferencia de despliegue es un forward pass. El sistema escala a
flotas grandes sin recalcular todo.

% ============================================================
\section{Modelo jerarquico bayesiano de flota multimarca}

El requerimiento de cualquier operador real: el sistema debe funcionar desde el
primer dia en cualquier marca, sin millones de datos por tipo de vehiculo.
La solucion matematica es \textbf{partial pooling bayesiano} en tres niveles:

\begin{mathbox}
$$\ln\eta_i = \underbrace{\alpha_{c,k}}_{\text{clase}} +
              \underbrace{\delta_{c,k,m}}_{\text{marca}} +
              \underbrace{\zeta_{c,k,m,j}}_{\text{modelo}} +
              \boldsymbol\gamma_c^\top\mathbf{x}_i$$
$$\alpha_{c,k} \sim \mathcal{N}(\mu_c,\tau_{\text{clase}}^2), \quad
  \delta_{c,k,m} \sim \mathcal{N}(0,\tau_{\text{marca}}^2), \quad
  \zeta_{c,k,m,j} \sim \mathcal{N}(0,\tau_{\text{modelo}}^2)$$
\end{mathbox}

El encogimiento opera en cascada: una marca con pocos datos toma fuerza prestada
de su clase, luego de la flota global. Las $\tau_\bullet$ (estimadas de los datos)
controlan cuanto pooling hay -- si las marcas se parecen, se agrupan fuerte;
si difieren, el modelo lo detecta. No hay parametros manuales.

Por que funciona entre clases tan distintas como camion pesado, vehiculo ligero
y motocicleta: \textbf{la fisica de falla es la misma} (Archard, Miner, Weibull,
Lundberg--Palmgren); solo cambian los parametros. Las features fisicas transfieren
entre clases aunque las marcas no se parezcan. Cuando una unidad tiene menos
sensores, la covariable faltante se marginaliza automaticamente:
$$p(T_i \mid \mathbf{x}_{i,\text{obs}}) =
  \int p(T_i \mid \mathbf{x}_{i,\text{obs}},x_{i,\text{mis}})
  \,p(x_{i,\text{mis}}\mid k)\,dx_{i,\text{mis}}$$

Una unidad con menos informacion produce un posterior predictivo mas ancho.
La incertidumbre se propaga correctamente a la RUL y al umbral de decision.

% ============================================================
\section{Decision optima: cuando intervenir}

\subsection{Intervalo optimo de reemplazo (teorema de renovacion-recompensa)}

Con costo preventivo $c_p$ y correctivo $c_f > c_p$, la tasa de costo a largo
plazo bajo politica de reemplazo por edad $T$:

\begin{mathbox}
$$g(T) = \frac{c_p\,R(T) + c_f\,[1-R(T)]}{\displaystyle\int_0^T R(t)\,dt}$$

La condicion de optimalidad $g'(T^*)=0$ reduce a:
$$h(T^*)\int_0^{T^*} R(t)\,dt - F(T^*) = \frac{c_p}{c_f - c_p}$$

El lado izquierdo crece con $T$ \textbf{solo si $h$ es creciente} ($\beta>1$,
desgaste). Por tanto existe $T^*<\infty$ si y solo si la pieza envejece.
Para $\beta=1$ (falla aleatoria), no existe optimo finito: el reemplazo
preventivo por calendario no ayuda.
\end{mathbox}

\subsection{Umbral de alerta sensible al costo}

El modelo entrega $p = P(\text{falla en el horizonte}\mid\mathbf{x})$.
La alerta optima es:

\begin{mathbox}
$$p > p^* = \frac{c_{\mathrm{FP}}}{c_{\mathrm{FP}} + c_{\mathrm{FN}}}$$

$p^*$ no es 0.5. En el reto publico de Scania APS (UCI, 2016),
$c_{\mathrm{FN}} = 500$, $c_{\mathrm{FP}} = 10$, de modo que $p^* \approx 2\%$:
se alerta con apenas 2\% de probabilidad de falla porque dejarla pasar
es 50 veces mas caro. La metrica correcta de entrenamiento no es exactitud:
es el costo monetario esperado.
\end{mathbox}

% ============================================================
\section{Cuantificacion del ahorro esperado}

Esta seccion construye el caso economico desde la matematica del modelo, no
desde numeros redondos de marketing. Se calculan tres horizontes.

\subsection{Modelo de costo total de una flota}

Sea $\mathcal{F}$ una flota de $N$ unidades. Para cada unidad $i$ y componente
$j$, definimos el costo total de mantenimiento en el horizonte $[0,T_H]$:

$$C_{\mathrm{total}} = \sum_{i=1}^{N}\sum_{j} \left[
  c_f^{(j)}\,E[N_f^{(ij)}] + c_p^{(j)}\,E[N_p^{(ij)}] + c_d\,E[D^{(ij)}]
\right]$$

donde $N_f^{(ij)}$ son fallas no anticipadas del componente $j$ en unidad $i$,
$N_p^{(ij)}$ son intervenciones preventivas, $D^{(ij)}$ son dias de tiempo muerto,
y $c_d$ es el costo diario de unidad varada.

\subsection{Relacion entre estrategia y $E[N_f]$}

Bajo mantenimiento reactivo puro, cada componente opera hasta la falla; el
numero esperado de fallas en $[0,T_H]$ es $E[N_f] = T_H / \mathrm{MTTF}$ (teoria
de renovacion). Bajo politica predictiva optima con umbral $R^*$, la intervencion
ocurre cuando $R(t)$ cae a $R^*$ antes de la falla; la probabilidad de falla
residual al momento de intervencion es $1-R^* \ll 1$.

El ahorro por componente por unidad en el horizonte $T_H$ es:

\begin{mathbox}
$$\Delta C^{(j)} = (c_f^{(j)} - c_p^{(j)})\cdot
  \underbrace{\left(\frac{T_H}{\mathrm{MTTF}^{(j)}} - \frac{T_H}{T^{*(j)}}\right)}
  _{\text{fallas evitadas}} \cdot \underbrace{(1 - \varepsilon)}_{\text{frac. anticipadas}}
  - c_p^{(j)}\cdot\underbrace{\frac{T_H}{T^{*(j)}}}_{\text{interv. preventivas}}
  + c_d \cdot \Delta D^{(j)}$$

donde $\varepsilon$ es la tasa de falsos negativos del modelo (fallas que el
sistema no anticipa), $T^{*(j)}$ es el intervalo optimo de la Parte~IX, y
$\Delta D^{(j)}$ es la reduccion de dias de tiempo muerto.
\end{mathbox}

\subsection{Parametros conservadores derivados de la literatura}

Los parametros siguientes son los minimos documentados por US DOE, McKinsey (2020)
y operadores de flota Clase 4--8. Se usa el rango conservador, no el optimista.

\begin{center}
\small
\begin{tabular}{@{}p{6.2cm}lll@{}}
\toprule
\textbf{Parametro} & \textbf{Conservador} & \textbf{Medio} & \textbf{Fuente} \\
\midrule
Reduccion de fallas no planeadas & 40\% & 70\% & US DOE \\
Reduccion del costo de reparacion & 25\% & 40\% & US DOE / McKinsey 2020 \\
Reduccion de tiempo muerto no planeado & 35\% & 50\% & US DOE / McKinsey 2020 \\
Multiplicador costo reactivo vs preventivo & $4.0\times$ & $4.8\times$ & US DOE \\
Extension de vida util del activo & 20\% & 40\% & McKinsey 2020 \\
Reduccion de inventario de refacciones & 20\% & 35\% & operadores flota \\
\bottomrule
\end{tabular}
\end{center}

\subsection{Ejercicio numerico: flota de referencia de 50 unidades}

Supongamos una flota de $N=50$ camiones de clase 6--8 en operacion en Mexico.
Parametros de entrada conservadores basados en datos de mercado mexicano (2025):

\begin{center}
\small
\begin{tabular}{@{}p{7.0cm}r@{}}
\toprule
\textbf{Parametro} & \textbf{Valor} \\
\midrule
Gasto anual en mantenimiento correctivo por unidad & MXN 85\,000 \\
Dias de tiempo muerto no planeado por unidad/ano & 6 dias \\
Costo de unidad varada por dia (ingreso no generado) & MXN 12\,000 \\
Costo promedio de reparacion de emergencia (DPF, motor) & MXN 48\,000 \\
Numero de fallas criticas en ruta por unidad/ano & 2.8 \\
Costo de intervencion preventiva programada (por evento) & MXN 9\,000 \\
\bottomrule
\end{tabular}
\end{center}

\medskip
\textbf{Linea base reactiva (situacion actual):}
\begin{align*}
C_{\mathrm{reactivo}} &= N \times \left[C_{\mathrm{mant}} + C_{\mathrm{varada}} + C_{\mathrm{emergencias}}\right] \\
  &= 50 \times \left[85{,}000 + 6 \times 12{,}000 + 2.8 \times 48{,}000\right] \\
  &= 50 \times \left[85{,}000 + 72{,}000 + 134{,}400\right] \\
  &= 50 \times 291{,}400 = \mathbf{MXN~14{,}570{,}000~\text{al ano}}
\end{align*}

\textbf{Con mantenimiento predictivo (escenario conservador, US DOE):}

Aplicando los minimos documentados: 40\% de reduccion en fallas no planeadas,
35\% de reduccion en tiempo muerto, 25\% de reduccion en costo de reparacion.
Se anade el costo de las intervenciones preventivas generadas por el sistema
($\approx 1.5$ por unidad/ano en el escenario conservador):

\begin{align*}
C_{\mathrm{predictivo}} &= N \times \Bigl[
    C_{\mathrm{mant}} \times 0.75
  + C_{\mathrm{varada}} \times 0.65
  + C_{\mathrm{emergencias}} \times 0.60
  + 1.5 \times c_p \Bigr] \\
  &= 50 \times \Bigl[63{,}750 + 46{,}800 + 80{,}640 + 13{,}500\Bigr] \\
  &= 50 \times 204{,}690 = \mathbf{MXN~10{,}234{,}500~\text{al ano}}
\end{align*}

\medskip
\begin{greenbox}
\begin{center}
\textbf{Ahorro anual estimado (escenario conservador): MXN 4{,}335{,}500}\\[4pt]
\textbf{Ahorro por unidad/ano: MXN 86{,}710 \;$\approx$\; USD 4{,}300}\\[4pt]
\textbf{Reduccion total de costo operativo: 30\%}
\end{center}
\vspace{2pt}
{\small Los escenarios medio y optimista documentados llevan el ahorro a MXN 5.8M--7.2M anuales
para la misma flota. La literatura reporta ROI de 10:1 en implementaciones maduras
(US DOE); el calculo anterior usa los parametros mas bajos.}
\end{greenbox}

\subsection{Valor presente neto del programa a 5 anos}

Sea $I_0$ el costo de implementacion del sistema (hardware de captura,
integracion, licencias). Con ahorro anual $\Delta C$ y tasa de descuento $r$:

\begin{mathbox}
$$\mathrm{VPN} = -I_0 + \sum_{k=1}^{5} \frac{\Delta C_k}{(1+r)^k}, \qquad
  \Delta C_k = \Delta C \times (1 + g)^{k-1}$$

donde $g$ es la tasa de mejora anual del modelo (los modelos mejoran al acumular
datos; la literatura reporta 30\%--40\% de incremento en ROI del ano 1 al ano 2).
\end{mathbox}

Para $I_0 = $ MXN~600\,000 (estimacion conservadora de hardware + integracion
para 50 unidades, $\approx$~MXN~12\,000/unidad), $\Delta C = $ MXN~4\,335\,500,
$r = 10\%$, $g = 15\%$:

$$\mathrm{VPN}_{5\text{a}} = -600{,}000 + \frac{4{,}335{,}500}{1.10} +
  \frac{4{,}985{,}825}{1.21} + \ldots \approx \mathbf{MXN~17.8~\text{millones}}$$

El \textbf{punto de equilibrio} se alcanza en el primer trimestre del ano 1.
Para una flota de 250 unidades, el mismo calculo escala linealmente a
MXN~21.7M de ahorro anual con $\mathrm{VPN}_{5\text{a}} \approx$ MXN~89M.

\subsection{Componentes adicionales del ahorro no incluidos en el calculo base}

\begin{itemize}
  \item \textbf{Extension de vida del activo.} Un camion de MXN~1.8M con 20\%
    de extension de vida (documentada por McKinsey) difiere MXN~360\,000 en
    costo de reposicion por unidad. Para 50 unidades con reposicion promedio
    en 7 anos: MXN~25.7M en valor diferido a valor presente.
  \item \textbf{Reduccion de inventario.} La capacidad de anticipar demanda de
    refacciones reduce el inventario de seguridad 20\%--35\%; para una flota
    con MXN~800\,000 inmovilizados en partes, libera MXN~160\,000--280\,000
    en capital de trabajo.
  \item \textbf{Reduccion de primas de seguro.} Aseguradoras en EE.UU. y Europa
    otorgan descuentos de 8\%--15\% a flotas con programas certificados de
    mantenimiento predictivo. En Mexico este mercado esta emergiendo.
  \item \textbf{Cumplimiento normativo.} Las NOMs aplicables al autotransporte
    federal (SICT) exigen condiciones fisico-mecanicas documentadas. El sistema
    genera ese registro automaticamente.
\end{itemize}

% ============================================================
\section{Por que este enfoque supera a los existentes}

Aqui la comparacion es tecnica, no comercial.

\subsection{Frente a sistemas basados en reglas y umbrales OEM}

Los sistemas actuales de telemetria (Samsara, Geotab, PeopleNet) leen el bus
CAN y convierten DTC en alertas de texto. Eso es CBM de primera generacion:
\textit{deteccion}, no \textit{prediccion}. No modelan la trayectoria de
degradacion; no cuantifican la incertidumbre; no calculan RUL; no optimizan
el momento de intervencion. La alerta llega cuando el codigo ya esta activo,
no dias antes.

\subsection{Frente a modelos puramente data-driven (LSTM, XGBoost)}

Un LSTM entrenado en datos de flota de Navistar no generaliza a Kenworth sin
reentrenamiento masivo. La razon es que los datos de telemetria cruda son
altamente especificos de la plataforma. Nuestro enfoque usa la fisica como
capa de abstraccion: la energia de frenado es la misma cantidad fisica
independientemente del OEM; el dano de Miner es independiente de la marca
del chasis. El feature engineering guiado por fisica transfiere entre marcas
y reduce la cantidad de datos necesaria en uno o dos ordenes de magnitud.

Ademas, un modelo data-driven puro no puede distinguir si una prediccion es
incierta porque el componente esta en zona de riesgo alto o porque simplemente
no hay datos de esa marca todavia. El modelo bayesiano lo sabe: el posterior
es mas ancho cuando hay menos datos, y esa incertidumbre se propaga
correctamente a la decision de intervencion.

\subsection{Frente a modelos de supervivencia sin fisica}

La supervivencia pura (Kaplan--Meier, Weibull sin covariables) da estimaciones
poblacionales pero no puede personalizar al individuo. Un Cox sin features
fisicas usa covariables de bajo poder predictivo (marca, ano) y no captura
la historia de carga real. La combinacion PoF~+~Cox~+~jerarquia bayesiana
es el enfoque state-of-the-art documentado en la literatura (Bellani et al.~2021;
Nascimento \& Viana, 2019; revision Theissler et al.~2021).

\subsection{El diferenciador central}

\begin{greenbox}
\textbf{Ninguna otra plataforma de tracking vehicular cierra el lazo
completo:} sensor $\to$ feature fisico (Archard, Miner, L10) $\to$ supervivencia
jerarquica bayesiana $\to$ RUL con intervalo de confianza $\to$ decision de costo
optimo $\to$ actualizacion continua en streaming con cada nuevo evento.

La mayoria opera hasta la tercera flecha. Nosotros operamos las cinco.
\end{greenbox}

\vspace{0.8cm}
\begin{center}
{\color{bordercol}\rule{0.35\textwidth}{0.5pt}}\\[0.35cm]
{\small\color{muted}Junio 2026 \quad|\quad Confidencial}
\end{center}

\end{document}
